\section{家户与神明之间：跨越两汉的保护关系}
\label{sec:han-households-religion-longue}

若只沿着诏令与战役阅读两汉，宗教、亲属和地方互助似乎总在政治之外。其实，
它们决定人们在征役、灾荒、迁徙和战争来临时向谁求助，也决定国家的命令会在
何处被解释、拖延或获得合作。家户不是一个封闭的私域：婚姻、继承、债务、祭祀、
佃作和收养把它连到乡里、市场、宗族与官府。神明、经典和治病仪式也不是彼此
隔绝的“信仰选项”，而是声望、书写、财物和保护可以交汇的场所。

\subsection{坟墓、祠祀与可见的亲属}

墓葬、画像石、碑刻和随葬文字让我们接触到汉人如何安排死者、祖先和家名，却
不能直接等同于全体居民的生活。能留下石刻和墓室的家庭本就拥有较多资源；没有
留下遗存的人并非没有亲属、祭祀或记忆，只是更难被今天看见。把考古材料同正史
的礼制与列传并读，能避免两个错误：一是把官修礼仪当作每一处乡里的统一实践，
二是把某一座墓的形制推作全国社会的平均面貌。

这种不均衡与政治编年密切相关。西汉迁徙关东豪族入关中、河西安置降附部众、东汉
战后重新编户，都使人们要在陌生地点重组家内劳力、婚姻和祭祀。官府需要这些
家户纳税、服役、听从司法；家户则可能通过亲属、同乡、雇佣和地方大户取得比
官方程序更快的保护。所谓豪强既能捐输、传递消息，也能使资源绕过直接征发，
正因为他们处在这种关系网的交叉处。

\subsection{知识如何成为庇护}

经学的师承、太学的空间和察举的评价，使文字与官职发生连接。白虎观和熹平石经
展示朝廷试图公开校定何种经典可以被援引；它们并不证明全国士人只有一种读法，
更不证明乡里社会只通过儒家语言行动。师友、家学、地方名望和捐献可以让不同
人进入官署，也能在官署拒绝他们时成为另一种声望来源。党锢之所以不只是几次
逮捕，正是因为禁令切断了部分地方评价通往中枢的路径。

楚王英案中的佛教供养字样、黄老方术、墓葬中的图像以及后来的道教材料，须各按
自己的年代和文本位置阅读。它们可以说明某些术语、物品或仪式已经被人使用，
不能给出信众数量、统一教团或私人信念的精确统计。材料的限度反而使问题更清楚：
当疾病、死亡、赋役和迁徙使普通家庭感到官府难以提供确定保护时，哪些仪式、
师徒和文字能把不安组织成可共享的关系？

\subsection{黄巾之后仍有许多不同的乡里}

太平道与黄巾起事使这个问题在一八四年成为朝廷危机。宗教语言可以把治病、互助、
灾异解释和末世期待接到行动上，但不能让所有参与者拥有同一动机或同一组织。
冀、豫、兖等地的道路、收成、乡里关系和既有武装各不相同；有人加入、有人回避、
有人被征发、有人借混乱改变与地主或官吏的关系。把黄巾只写成“农民战争”或
“邪教叛乱”，都会把这些差异压平。

镇压也没有让这些保护问题消失。州郡需要自保，地方精英需要处理流民和军粮，
军人需要寻找供给，宗族与宾客网络因而获得更大作用。东汉末年的区域竞争并非
社会关系突然取代国家，而是国家原有的簿籍、道路、荐举和救济接口被不同集团
分别占有。二二〇年的禅代记录了一个皇统的结束，不能结束家庭如何寻找土地、
祭祀、工作和保护的历史；魏、蜀、吴都会在这些已经改变的关系上继续组织国家。

\begin{textwindow}[读社会史，也读它的缺页]
正史最善于保存诏令、战事和显贵人物；简牍、墓葬、碑刻、钱币与宗教文本把不同
地点的劳动、书写和仪式局部带回视野。没有一种材料能替沉默者说完故事。把它们
放入编年叙事，不是给帝王史加一层装饰，而是追问一项命令为何在不同地方遇到
不同的家户、道路和保护关系。
\end{textwindow}
